\documentclass[11pt]{article}

\usepackage[margin=1in]{geometry}
\usepackage{amsmath,amssymb,amsthm}
\usepackage{booktabs}
\usepackage{graphicx}
\usepackage{hyperref}
\usepackage{cleveref}
\usepackage{algorithm}
\usepackage{algpseudocode}
\usepackage{xcolor}

\newtheorem{definition}{Definition}

\newcommand{\E}{\mathcal{E}}
\newcommand{\Qemp}{Q_{\mathrm{emp}}}
\newcommand{\argmax}{\operatorname{arg\,max}}
\newcommand{\nats}{\,\mathrm{nats}}

\title{Empowerment Outer-Policy Experiments:\\
  Separating Inner Valuation from Outer Control}
\author{Scott Viteri}
\date{April 2026}

\begin{document}
\maketitle

\begin{abstract}
Empowerment---the channel capacity of an agent's action-to-future-state
channel---is a widely studied intrinsic motivation signal.  However, prior
work often conflates the \emph{empowerment objective} (which states are
valued) with the \emph{outer control rule} (how actions are selected given
empowerment estimates).  This repository provides a controlled comparison:
we hold empowerment computation fixed and vary only the outer policy.
We compare greedy (argmax), softmax (Boltzmann), and $\epsilon$-greedy
outer policies on exact tabular empowerment in a corridor gridworld.
Greedy control collapses into stereotyped local loops, while moderate-temperature
softmax control escapes local empowerment maxima and reaches higher-empowerment
regions at the cost of reduced instantaneous empowerment.
\end{abstract}

%----------------------------------------------------------------------
\section{Introduction}
\label{sec:intro}
%----------------------------------------------------------------------

Empowerment \cite{klyubin2005empowerment} provides an agent-centric,
task-independent measure of how much control an agent has over its
future.  It has been used to explain and generate life-like behaviors
such as obstacle avoidance, resource acquisition, and survival
\cite{jung2011empowerment, mohamed2015variational, ramirez2024mop}.

A subtle but important design choice in empowerment-driven agents is the
\emph{outer policy}: given empowerment estimates for all reachable
states, how should the agent select its next action?  The simplest
choice is greedy (argmax) over successor-state empowerment.  But
stochastic alternatives---softmax or $\epsilon$-greedy---can produce
qualitatively different behavior even when the underlying empowerment
values are identical.

This repository isolates this variable.  The central experimental
principle is:

\begin{quote}
\textbf{Fix empowerment.  Vary only the outer policy.  Measure behavior broadly.}
\end{quote}

%----------------------------------------------------------------------
\section{Mathematical Framework}
\label{sec:math}
%----------------------------------------------------------------------

\subsection{Empowerment}

Consider a discrete environment with state space $\mathcal{S}$ and
action space $\mathcal{A}(s)$ available at each state $s$.  The
transition dynamics are $P(s' \mid s, a)$.

\begin{definition}[n-step empowerment]
For state $s$ and horizon $n$, the \emph{empowerment} is
\begin{equation}
  \E_n(s) \;=\; \max_{q(a_{0:n-1} \mid s)}\; I(A_{0:n-1};\, S_n \mid s)
  \label{eq:empowerment}
\end{equation}
where $I(\cdot;\cdot)$ is mutual information, $A_{0:n-1}$ is a
length-$n$ action sequence drawn from the input distribution
$q(a_{0:n-1} \mid s)$, and $S_n$ is the resulting final state under
the environment dynamics.
\end{definition}

Empowerment is the \emph{channel capacity} of the $n$-step
action-to-state channel induced by the dynamics at $s$.  It measures,
in nats, the maximum number of distinguishable futures the agent can
steer toward.

\subsection{The Action-to-State Channel}

For a given state $s$ and horizon $n$, we construct the channel matrix
$\mathbf{C} \in \mathbb{R}^{|\mathcal{X}| \times |\mathcal{S}|}$
where $\mathcal{X} = \mathcal{A}(s)^n$ is the set of all length-$n$
action sequences.  Each entry is:
\begin{equation}
  C_{x,s'} = P(S_n = s' \mid s,\, a_{0:n-1} = x)
\end{equation}
computed by rolling out the transition probabilities:
\begin{equation}
  P(S_n = s' \mid s, x) \;=\;
    \sum_{s_1, \ldots, s_{n-1}}
    \prod_{t=0}^{n-1} P(s_{t+1} \mid s_t, a_t)
\end{equation}
where $s_0 = s$ and $x = (a_0, \ldots, a_{n-1})$.  In deterministic
environments, each row of $\mathbf{C}$ is a one-hot vector.

\subsection{Blahut--Arimoto Algorithm}

Channel capacity is computed via the Blahut--Arimoto (BA) algorithm
\cite{blahut1972computation}.  Starting from a uniform input
distribution $q^{(0)}$, BA iterates:

\begin{enumerate}
  \item Compute the output marginal:
    $r^{(t)}(s') = \sum_x q^{(t)}(x)\, C_{x,s'}$
  \item Compute the per-input KL term:
    $\phi^{(t)}(x) = \exp\!\Big(\sum_{s'} C_{x,s'} \ln
    \frac{C_{x,s'}}{r^{(t)}(s')}\Big)$
  \item Update input distribution:
    $q^{(t+1)}(x) \propto q^{(t)}(x)\, \phi^{(t)}(x)$
\end{enumerate}

Convergence is monitored via the gap between the lower bound
$C_L = \ln\!\big(\sum_x q(x)\phi(x)\big)$ and upper bound
$C_U = \ln\!\big(\max_x \phi(x)\big)$.  The algorithm terminates when
$C_U - C_L < \epsilon$ (we use $\epsilon = 10^{-8}$).

\subsection{Empowerment-Derived Action Scores}

Given a precomputed empowerment table
$\{\E_n(s)\}_{s \in \mathcal{S}}$, we define the one-step action
score:
\begin{equation}
  \Qemp(s, a) \;=\; \sum_{s'} P(s' \mid s, a)\, \E_n(s')
  \label{eq:qemp}
\end{equation}
In deterministic environments this simplifies to
$\Qemp(s,a) = \E_n(f(s,a))$ where $f$ is the transition function.

\textbf{This score is the same for all outer policies}---the only
thing that differs is how $\Qemp$ is converted into action
probabilities.

%----------------------------------------------------------------------
\section{Outer Policies}
\label{sec:policies}
%----------------------------------------------------------------------

\subsection{Greedy (Argmax)}

The greedy policy selects uniformly among maximizers:
\begin{equation}
  \pi_{\mathrm{greedy}}(a \mid s) \;=\;
    \frac{\mathbf{1}[a \in \argmax_b \Qemp(s,b)]}{|\argmax_b \Qemp(s,b)|}
  \label{eq:greedy}
\end{equation}

\subsection{Softmax (Boltzmann)}

The softmax policy assigns probability proportional to exponentiated
scores:
\begin{equation}
  \pi_\tau(a \mid s) \;=\;
    \frac{\exp\!\big(\Qemp(s,a) / \tau\big)}
         {\sum_b \exp\!\big(\Qemp(s,b) / \tau\big)}
  \label{eq:softmax}
\end{equation}
where $\tau > 0$ is the temperature parameter.  Key limits:
\begin{itemize}
  \item $\tau \to 0$: recovers the greedy policy
  \item $\tau \to \infty$: approaches the uniform distribution over actions
\end{itemize}

For numerical stability, we subtract $\max_b \Qemp(s,b)$ from all
scores before exponentiation.

\subsection{Epsilon-Greedy}

The $\epsilon$-greedy policy serves as a control to distinguish
``stochasticity helps'' from ``Boltzmann-shaped stochasticity helps'':
\begin{equation}
  \pi_\epsilon(a \mid s) \;=\;
  \begin{cases}
    \dfrac{1 - \epsilon}{|\argmax_b \Qemp(s,b)|} + \dfrac{\epsilon}{|\mathcal{A}(s)|}
      & a \in \argmax_b \Qemp(s,b) \\[6pt]
    \dfrac{\epsilon}{|\mathcal{A}(s)|}
      & \text{otherwise}
  \end{cases}
  \label{eq:epsgreedy}
\end{equation}

%----------------------------------------------------------------------
\section{Corridor Environment}
\label{sec:corridor}
%----------------------------------------------------------------------

The corridor environment is a deterministic gridworld designed to
create the simplest possible test of local-maximum escape.  It
consists of:

\begin{itemize}
  \item A \textbf{local basin}: a $7 \times 4$ open room (28 states)
    with moderate empowerment, containing the start state.
  \item A \textbf{narrow corridor}: 3 cells wide at a single row,
    connecting the two basins.  These states have reduced empowerment
    because fewer future states are reachable within the horizon.
  \item A \textbf{far basin}: a $7 \times 7$ open room (49 states)
    with higher empowerment due to its larger spatial extent.
\end{itemize}

The agent has four actions (up, down, left, right) at every state.
Moving into a wall results in staying in place.  There are no
absorbing states.

\medskip

At horizon $n = 3$, the mean empowerment values are:
\begin{itemize}
  \item Local basin: $2.631\nats$
  \item Far basin: $2.693\nats$
\end{itemize}

The far basin has higher empowerment because its larger area provides
more reachable states within any given horizon.  However, to reach it,
the agent must traverse corridor states with \emph{lower} empowerment
than the local basin---a greedy agent will not voluntarily enter them.

%----------------------------------------------------------------------
\section{Experimental Setup}
\label{sec:setup}
%----------------------------------------------------------------------

We precompute the exact $n$-step empowerment table for all 80 states
using BA with horizon $n = 3$.  We then run rollouts of length
$T = 200$ steps for each outer policy across 100 random seeds.

\paragraph{Independent variables.}
\begin{itemize}
  \item Outer policy: greedy, softmax ($\tau \in \{0.01, 0.03, 0.1, 0.3, 1.0, 3.0\}$),
    $\epsilon$-greedy ($\epsilon \in \{0.1, 0.3\}$)
\end{itemize}

\paragraph{Dependent variables.}
\begin{itemize}
  \item Fraction of episodes reaching the far basin
  \item Mean state visitation entropy (Shannon, in nats)
  \item Mean number of distinct states visited
  \item Mean empowerment of visited states
\end{itemize}

%----------------------------------------------------------------------
\section{Results}
\label{sec:results}
%----------------------------------------------------------------------

\begin{table}[h]
\centering
\caption{Corridor experiment results ($n=3$, $T=200$, 100 seeds).}
\label{tab:results}
\begin{tabular}{@{}lcccc@{}}
\toprule
Policy & Far Basin \% & Entropy & Distinct States & Mean Emp. \\
\midrule
Greedy              & 0.00 & 1.301 &  6.0 & 2.861 \\
Softmax $\tau=0.01$ & 0.00 & 1.311 &  6.3 & 2.860 \\
Softmax $\tau=0.03$ & 0.00 & 1.646 &  9.1 & 2.853 \\
Softmax $\tau=0.1$  & 0.00 & 2.357 & 16.7 & 2.820 \\
Softmax $\tau=0.3$  & 0.10 & 2.961 & 26.0 & 2.751 \\
Softmax $\tau=1.0$  & 0.29 & 3.269 & 33.7 & 2.682 \\
Softmax $\tau=3.0$  & 0.33 & 3.255 & 33.6 & 2.652 \\
\midrule
$\epsilon$-greedy $\epsilon=0.1$ & 0.02 & 1.516 & 10.3 & 2.852 \\
$\epsilon$-greedy $\epsilon=0.3$ & 0.25 & 2.027 & 16.9 & 2.831 \\
\bottomrule
\end{tabular}
\end{table}

The results in \Cref{tab:results} support the core hypotheses:

\paragraph{H1: Greedy empowerment is behaviorally narrow.}
The greedy agent visits only 6 distinct states out of 80 and never
reaches the far basin.  It becomes trapped in a small loop near the
highest-empowerment states of the local basin.

\paragraph{H2: Softmax improves basin escape.}
At $\tau = 0.3$, 10\% of episodes reach the far basin; at
$\tau = 1.0$, this rises to 29\%.  The softmax policy's willingness
to occasionally accept lower-empowerment actions allows it to traverse
the corridor.

\paragraph{H3: Diversity--empowerment trade-off.}
Visitation entropy increases monotonically from $1.30\nats$ (greedy)
to $3.27\nats$ ($\tau = 1.0$), while mean empowerment of visited
states decreases from $2.86\nats$ to $2.68\nats$.  Softmax explores
more broadly but spends less time at peak-empowerment states.

\paragraph{H4: Diminishing returns at high temperature.}
Between $\tau = 1.0$ and $\tau = 3.0$, the far-basin reach rate
plateaus (29\% $\to$ 33\%) while entropy saturates and mean
empowerment continues to decline.  Very high temperatures approach
uniform random behavior.

\paragraph{Boltzmann vs.\ uniform stochasticity.}
At comparable diversity levels, $\epsilon$-greedy with $\epsilon = 0.3$
reaches the far basin 25\% of the time but visits only 16.9 distinct
states, whereas softmax at $\tau = 1.0$ reaches 29\% and visits 33.7
distinct states.  The Boltzmann weighting distributes exploration
probability more effectively than uniform noise.

%----------------------------------------------------------------------
\section{Repository Structure}
\label{sec:repo}
%----------------------------------------------------------------------

The implementation follows a strict separation between inner
empowerment computation and outer policy selection.

\begin{description}
  \item[\texttt{src/envs/}] Environment interface (\texttt{DiscreteEnv}
    base class) and the corridor gridworld.
  \item[\texttt{src/empowerment/}] Exact empowerment computation:
    \texttt{nstep\_model.py} builds the $n$-step channel matrix,
    \texttt{blahut\_arimoto.py} computes channel capacity, and
    \texttt{exact.py} ties them together.  \texttt{score.py} computes
    $\Qemp(s,a)$.
  \item[\texttt{src/policies/}] Outer policies (greedy, softmax,
    $\epsilon$-greedy), all consuming the same $\Qemp$ scores.
  \item[\texttt{src/rollout/}] Seeded simulation loop.
  \item[\texttt{src/metrics/}] Visitation entropy, distinct states,
    basin-reaching fraction.
  \item[\texttt{scripts/}] Experiment runner with temperature sweep
    and summary table output.
  \item[\texttt{tests/}] Unit tests for BA convergence, environment
    dynamics, and policy properties (17 tests).
\end{description}

%----------------------------------------------------------------------
\section{Discussion}
\label{sec:discussion}
%----------------------------------------------------------------------

These results demonstrate that behavioral conclusions about
empowerment-driven agents can depend substantially on the outer
control rule.  The greedy agent's stereotyped looping is not a
property of empowerment \emph{per se}---it is a property of greedy
action selection applied to empowerment scores.  A softmax agent using
the \emph{identical} empowerment values exhibits qualitatively
different behavior: broader exploration, basin escape, and higher
trajectory diversity.

This distinction matters for interpreting prior work.  When a paper
reports that ``empowerment-driven agents exhibit behavior X,'' the
outer policy is a confound.  Our framework makes it possible to ask:
is X caused by what empowerment values, or by how the agent acts on
those values?

\paragraph{Limitations.}
The current results use a single environment with a small empowerment
gap between basins ($\Delta \E \approx 0.06\nats$).  Larger
environments, longer horizons, and environments with survival pressure
(e.g., four-rooms with energy) would provide stronger tests.
The softmax temperature is also environment-specific; score
normalization (z-score mode) would improve cross-environment
comparability.

\paragraph{Future work.}
Planned extensions include four-rooms with energy constraints, a noisy
room variant, discretized cartpole, and horizon sweep ablations.

%----------------------------------------------------------------------
\bibliographystyle{plain}
\begin{thebibliography}{9}

\bibitem{klyubin2005empowerment}
A.~S. Klyubin, D.~Polani, and C.~L. Nehaniv.
\newblock Empowerment: A universal agent-centric measure of control.
\newblock In \emph{IEEE Congress on Evolutionary Computation}, 2005.

\bibitem{jung2011empowerment}
T.~Jung, D.~Polani, and P.~Stone.
\newblock Empowerment for continuous agent--environment systems.
\newblock \emph{Adaptive Behavior}, 19(1):16--39, 2011.

\bibitem{mohamed2015variational}
S.~Mohamed and D.~J. Rezende.
\newblock Variational information maximisation for intrinsically motivated
  reinforcement learning.
\newblock In \emph{NeurIPS}, 2015.

\bibitem{ramirez2024mop}
J.~Ram\'irez-Ruiz, C.~W.~J. Geurts, and R.~Moreno-Bote.
\newblock Complex behavior from intrinsic motivation to occupy future
  action-state path space.
\newblock \emph{Nature Communications}, 2024.

\bibitem{blahut1972computation}
R.~Blahut.
\newblock Computation of channel capacity and rate-distortion functions.
\newblock \emph{IEEE Transactions on Information Theory}, 18(4):460--473, 1972.

\end{thebibliography}

\end{document}
